% GENERATED FILE. Edit canonical lessons, then run npm run generate:books.
% canonical-source-hash: fnv1a64:47bc1551d7edc210
% canonical-lessons: BN-C363-somosya, BN-C363-upay, BN-C363-bikolpo, BN-C363-uttor, BN-C363-utre-jaoya

\chapter{Problems and Answers}
\label{ch:bn-problems-and-answers}
\hlchaptermodality{363}

\begin{chapteropening}
\textbf{By the end of this chapter:} \emph{I can name a problem and suggest ways out of it.}

Complete the last lesson of chapter 363: I can name a problem and suggest ways out of it.
\end{chapteropening}

\section[\texorpdfstring{\bn{সমস্যা} (sômôsyā)}{সমস্যা (sômôsyā)}]{\bn{সমস্যা} (sômôsyā) — a problem}

\label{lesson:BN-C363-somosya}



\begin{quote}
Before the new one: say the Bengali for to sense, then the Bengali for clear.
\end{quote}

\subsection*{You'll want to know: \bn{সমস্যা}}
\textbf{\bn{সমস্যা}} — \emph{sômôsyā} — \textquotedblleft{}a problem\textquotedblright{}.

\begin{quote}
\bn{তোমার} \bn{কোনো} \bn{সমস্যা} \bn{হলে} \bn{আমাকে} \bn{বলো।}
\end{quote}

\begin{quote}
\emph{tomār kono sômôsyā hôle āmāke bôlo.}
\end{quote}

\begin{quote}
If you have any problem, tell me.
\end{quote}

\begin{cousinweb}
\textbf{\bn{সমস্যা}} (\emph{sômôsyā}) is Sanskrit \emph{samasyā}, a part of a verse given to someone else to complete, a kind of puzzle (Monier-Williams). Bengali widened the puzzle to any problem; that widening is not in the dictionary.
\end{cousinweb}

\begin{grammarlens}[title={What you've built}]
One more word for giving reasons and opinions.
\end{grammarlens}

\subsection*{Your turn}
\begin{itemize}
\raggedright
  \item \emph{Say it:} \emph{sômôsyā}
  \item \emph{Say it:} \emph{sômôsyā}, once more
  \item \emph{Say it:} \emph{tomār kono sômôsyā hôle āmāke bôlo.} — If you have any problem, tell me.
  \item \emph{Recall:} say the Bengali for to sense, then the Bengali for clear, then say \emph{sômôsyā} again
\end{itemize}

\begin{tcolorbox}[breakable,colback=teal!4,colframe=teal!35!black,title={Before you move on}]
What does \bn{সমস্যা} mean? (\textquotedblleft{}A problem\textquotedblright{}.) Say it once more.
\end{tcolorbox}
\section[\texorpdfstring{\bn{উপায়} (upāy)}{উপায় (upāy)}]{\bn{উপায়} (upāy) — a way; a means}

\label{lesson:BN-C363-upay}



\begin{quote}
Before the new one: say the Bengali for clear, then the Bengali for a problem.
\end{quote}

\subsection*{You'll want to know: \bn{উপায়}}
\textbf{\bn{উপায়}} — \emph{upāy} — \textquotedblleft{}a way; a means\textquotedblright{}.

\begin{quote}
\bn{আর} \bn{কোনো} \bn{উপায়} \bn{নেই।}
\end{quote}

\begin{quote}
\emph{ār kono upāy nei.}
\end{quote}

\begin{quote}
There is no other way.
\end{quote}

\begin{cousinweb}
\textbf{\bn{উপায়}} (\emph{upāy}) is Sanskrit \emph{upāya}, 'coming near', a means of reaching an end (Monier-Williams).
\end{cousinweb}

\begin{grammarlens}[title={What you've built}]
One more word for giving reasons and opinions.
\end{grammarlens}

\subsection*{Your turn}
\begin{itemize}
\raggedright
  \item \emph{Say it:} \emph{upāy}
  \item \emph{Say it:} \emph{upāy}, once more
  \item \emph{Say it:} \emph{ār kono upāy nei.} — There is no other way.
  \item \emph{Recall:} say the Bengali for clear, then the Bengali for a problem, then say \emph{upāy} again
\end{itemize}

\begin{tcolorbox}[breakable,colback=teal!4,colframe=teal!35!black,title={Before you move on}]
What does \bn{উপায়} mean? (\textquotedblleft{}A way; a means\textquotedblright{}.) Say it once more.
\end{tcolorbox}
\section[\texorpdfstring{\bn{বিকল্প} (bikôlpô)}{বিকল্প (bikôlpô)}]{\bn{বিকল্প} (bikôlpô) — an alternative; an option}

\label{lesson:BN-C363-bikolpo}



\begin{quote}
Before the new one: say the Bengali for a problem, then the Bengali for a means.
\end{quote}

\subsection*{You'll want to know: \bn{বিকল্প}}
\textbf{\bn{বিকল্প}} — \emph{bikôlpô} — \textquotedblleft{}an alternative; an option\textquotedblright{}.

\begin{quote}
\bn{সাইকেল} \bn{গাড়ির} \bn{একটা} \bn{ভালো} \bn{বিকল্প।}
\end{quote}

\begin{quote}
\emph{sāikel gāṛir ekṭā bhālo bikôlpô.}
\end{quote}

\begin{quote}
A bicycle is a good alternative to a car.
\end{quote}

\begin{cousinweb}
\textbf{\bn{বিকল্প}} (\emph{bikôlpô}) is Sanskrit \emph{vikalpa}, an alternative or option (Monier-Williams). What it replaces takes -\bn{এর}.
\end{cousinweb}

\begin{grammarlens}[title={What you've built}]
One more word for giving reasons and opinions.
\end{grammarlens}

\subsection*{Your turn}
\begin{itemize}
\raggedright
  \item \emph{Say it:} \emph{bikôlpô}
  \item \emph{Say it:} \emph{bikôlpô}, once more
  \item \emph{Say it:} \emph{sāikel gāṛir ekṭā bhālo bikôlpô.} — A bicycle is a good alternative to a car.
  \item \emph{Recall:} say the Bengali for a problem, then the Bengali for a means, then say \emph{bikôlpô} again
\end{itemize}

\begin{tcolorbox}[breakable,colback=teal!4,colframe=teal!35!black,title={Before you move on}]
What does \bn{বিকল্প} mean? (\textquotedblleft{}An alternative; an option\textquotedblright{}.) Say it once more.
\end{tcolorbox}
\section[\texorpdfstring{\bn{উত্তর} (uttôr)}{উত্তর (uttôr)}]{\bn{উত্তর} (uttôr) — an answer; a reply}

\label{lesson:BN-C363-uttor}



\begin{quote}
Before the new one: say the Bengali for a means, then the Bengali for an alternative.
\end{quote}

\subsection*{You'll want to know: \bn{উত্তর}}
\textbf{\bn{উত্তর}} — \emph{uttôr} — \textquotedblleft{}an answer; a reply\textquotedblright{}.

\begin{quote}
\bn{আমার} \bn{প্রশ্নের} \bn{উত্তর} \bn{দাও।}
\end{quote}

\begin{quote}
\emph{āmār prôshner uttôr dāo.}
\end{quote}

\begin{quote}
Answer my question.
\end{quote}

\begin{cousinweb}
\textbf{\bn{উত্তর}} (\emph{uttôr}) is Sanskrit \emph{uttara}, the later or the upper one, and so a reply (Monier-Williams). The same word also means north. It is given with \bn{দেওয়া}.
\end{cousinweb}

\begin{grammarlens}[title={What you've built}]
One more word for giving reasons and opinions.
\end{grammarlens}

\subsection*{Your turn}
\begin{itemize}
\raggedright
  \item \emph{Say it:} \emph{uttôr}
  \item \emph{Say it:} \emph{uttôr}, once more
  \item \emph{Say it:} \emph{āmār prôshner uttôr dāo.} — Answer my question.
  \item \emph{Recall:} say the Bengali for a means, then the Bengali for an alternative, then say \emph{uttôr} again
\end{itemize}

\begin{tcolorbox}[breakable,colback=teal!4,colframe=teal!35!black,title={Before you move on}]
What does \bn{উত্তর} mean? (\textquotedblleft{}An answer; a reply\textquotedblright{}.) Say it once more.
\end{tcolorbox}
\section[\texorpdfstring{\bn{উতরে} \bn{যাওয়া} (utre jāoyā)}{উতরে যাওয়া (utre jāoyā)}]{\bn{উতরে} \bn{যাওয়া} (utre jāoyā) — to come through; to get through (a test)}

\label{lesson:BN-C363-utre-jaoya}



\begin{quote}
Before the new one: say the Bengali for an alternative, then the Bengali for an answer.
\end{quote}

\subsection*{You'll want to know: \bn{উতরে} \bn{যাওয়া}}
\textbf{\bn{উতরে} \bn{যাওয়া}} — \emph{utre jāoyā} — \textquotedblleft{}to come through; to get through (a test)\textquotedblright{}.

\begin{quote}
\bn{অনেক} \bn{ঝামেলার} \bn{পরে} \bn{আমরা} \bn{উতরে} \bn{গেলাম।}
\end{quote}

\begin{quote}
\emph{ônek jhāmelār pôre āmrā utre gelām.}
\end{quote}

\begin{quote}
After a lot of trouble, we came through.
\end{quote}

Past (he/she …): \bn{সেই} \bn{পরিকল্পনাটা} \bn{উতরে} \bn{গেল}, \emph{sei pôrikôlpônāṭā utre gelô}.

\begin{cousinweb}
\textbf{\bn{উতরে} \bn{যাওয়া}} (\emph{utre jāoyā}) is the -\bn{এ} form of \bn{উতরা} 'to cross, to get over' (CDIAL 1770 lists Bengali utarā) with \bn{যাওয়া}: to get across a difficulty to the other side. The sense 'to pass a test' is a Bengali use.
\end{cousinweb}

\begin{grammarlens}[title={What you've built}]
One more word for giving reasons and opinions.
\end{grammarlens}

\subsection*{Your turn}
\begin{itemize}
\raggedright
  \item \emph{Say it:} \emph{utre jāoyā}
  \item \emph{Say it:} \emph{utre jāoyā}, once more
  \item \emph{Say it:} \emph{ônek jhāmelār pôre āmrā utre gelām.} — After a lot of trouble, we came through.
  \item \emph{Recall:} say the Bengali for an alternative, then the Bengali for an answer, then say \emph{utre jāoyā} again
\end{itemize}

\begin{tcolorbox}[breakable,colback=teal!4,colframe=teal!35!black,title={Before you move on}]
What does \bn{উতরে} \bn{যাওয়া} mean? (\textquotedblleft{}To come through; to get through (a test)\textquotedblright{}.) Say it once more.
\end{tcolorbox}
